\section{Coding Theory}\label{sec:coding_theory}

This section records the coding-theoretic interface used by the FRI, STIR, and WHIR
formalizations. The Lean development is intentionally split between a small reusable core
(``distance'', ``linear code'', ``list decoding'', and Reed--Solomon code APIs) and a larger
proximity-gap/list-decoding layer. The proof-system chapters use the core definitions directly
and consume the deeper proximity-gap theorems through explicit named hypotheses when the
corresponding mathematics is still open in ArkLib.

\subsection{Core codes and distances}\label{sec:coding_core}

\begin{definition}[Code distance]\label{def:code_distance}
    \lean{Code.dist,Code.relHammingDist}
    The Hamming distance $\Delta_0(f,g)$ and relative Hamming distance $\Delta(f,g)$ count,
    respectively, the number and fraction of coordinates on which two words disagree.
\end{definition}

\begin{definition}[Distance from a code]\label{def:distance_from_code}
    \lean{Code.distFromCode,Code.relDistFromCode}
    \uses{def:code_distance}
    The distance from a word $f$ to a code $\code$ is the minimum distance from $f$ to a
    codeword of $\code$, in either absolute or relative form.
\end{definition}

\begin{definition}[Abstract code]\label{def:code}
    \lean{ListDecodable.Code}
    A code over an index type $\iota$ and alphabet $F$ is represented as a set of words
    $\iota \to F$.
\end{definition}

\begin{definition}[Linear code]\label{def:linear_code}
    \lean{LinearCode}
    \uses{def:code}
    Linear codes are submodules of the word module. This is the representation used by
    Reed--Solomon codes and by the correlated-agreement predicates.
\end{definition}

\begin{definition}[Generator and check matrices]\label{def:generator_matrix}
    \lean{LinearCode.fromRowGenMat,LinearCode.fromColGenMat,LinearCode.byCheckMatrix}
    \uses{def:linear_code}
    Linear codes can be constructed from row or column generator matrices, and from parity-check
    matrices.
\end{definition}

\begin{definition}[Interleaved code]\label{def:interleaved_code}
    \lean{interleavedCodeSet,CodeInterleavable}
    \uses{def:linear_code}
    The interleaved code of a base code records stacks of codewords sharing the same coordinate
    set. It is the ambient object behind joint agreement and list-recovery style arguments.
\end{definition}

\subsection{Reed--Solomon codes}\label{sec:reed_solomon_codes}

\begin{definition}[Reed--Solomon code]\label{def:reed_solomon_code}
    \lean{ReedSolomon.evalOnPoints,ReedSolomon.code,ReedSolomon.genMatrix,ReedSolomon.checkMatrix}
    \uses{def:linear_code}
    For an injective evaluation map $\phi : \iota \hookrightarrow \field$ and a degree bound
    $d$, $\rscode[\field,\phi,d]$ is the image of polynomials of degree $<d$ under evaluation on
    $\phi(\iota)$. The library also provides the Vandermonde generator matrix and the standard
    check matrix for finite domains.
\end{definition}

\begin{lemma}[Codeword witnesses and transport]\label{lem:rs_codeword_api}
    \lean{ReedSolomon.mem_code_of_polynomial_of_degree_lt_of_eval,ReedSolomon.mem_code_iff_exists_polynomial,ReedSolomon.codeword_equiv_of_eval_eq}
    \uses{def:reed_solomon_code}
    A word is in a Reed--Solomon code when it is the evaluation of a polynomial below the degree
    bound. Codewords can also be transported across equivalent evaluation domains when the
    embeddings select the same field points.
\end{lemma}

\begin{definition}[Smooth and constrained Reed--Solomon codes]\label{def:smooth_rs_code}
    \lean{smoothCode,constrainedCode,multiConstrainedCode}
    \uses{def:reed_solomon_code}
    STIR and WHIR use smooth Reed--Solomon domains and constrained variants. The constrained
    code definitions add one or more polynomial constraints to the base smooth Reed--Solomon
    relation.
\end{definition}

\subsection{List decoding}\label{sec:list_decoding}

\begin{definition}[Hamming balls and close codewords]\label{def:list_close_codewords}
    \lean{ListDecodable.hammingBall,ListDecodable.relHammingBall,ListDecodable.closeCodewords,ListDecodable.closeCodewordsRel}
    \uses{def:distance_from_code,def:code}
    The list around a received word is the set of codewords whose Hamming or relative Hamming
    distance from that word is at most the chosen radius.
\end{definition}

\begin{definition}[List decodability]\label{def:list_decodable}
    \lean{ListDecodable.listDecodable,ListDecodable.uniqueDecodable}
    \uses{def:list_close_codewords}
    A code is $(\delta,\ell)$-list decodable when every relative Hamming ball of radius
    $\delta$ contains at most $\ell$ codewords. Unique decodability is the special case
    $\ell = 1$.
\end{definition}

\begin{definition}[Maximized list size]\label{def:lambda_list_size}
    \lean{ListDecodable.Lambda,ListDecodable.closeCodewordsRelFinset}
    \uses{def:list_decodable}
    The value $\Lambda(\code,\delta)$ is the maximum list size over all received words. Finset
    wrappers connect this maximized definition to finite-cardinality proofs.
\end{definition}

\begin{lemma}[Basic list-size monotonicity]\label{lem:list_size_basics}
    \lean{ListDecodable.closeCodewordsRel_subset_of_le,ListDecodable.Lambda_mono,ListDecodable.Lambda_le_natCast_of_forall_ncard_le}
    \uses{def:lambda_list_size}
    Close-codeword lists and the maximized list size are monotone in the decoding radius, and
    pointwise finite list-size bounds package into a bound on $\Lambda(\code,\delta)$.
\end{lemma}

\subsection{Proximity gaps and correlated agreement}\label{sec:proximity_gap_defs}

\begin{definition}[Proximity measure and proximity gap]\label{def:proximity_gap}
    \lean{ProximityGap.proximityMeasure,ProximityGap.proximityGap,ProximityGap.δ_ε_proximityGap}
    \uses{def:distance_from_code}
    The proximity-gap predicates express a dichotomy: a sampled family is either always
    $\delta$-close to the target code or is $\delta$-close with probability at most $\epsilon$.
    This is the generic coding-theory statement that FRI, STIR, and WHIR specialize to
    Reed--Solomon codes.
\end{definition}

\begin{definition}[Correlated agreement]\label{def:correlated_agreement}
    \lean{ProximityGap.δ_ε_correlatedAgreementAffineLines,ProximityGap.δ_ε_correlatedAgreementCurves,ProximityGap.δ_ε_correlatedAgreementAffineSpaces}
    \uses{def:proximity_gap,def:interleaved_code}
    Correlated agreement turns a high probability of random linear, curve, or affine-space
    combinations being close to a code into a common agreement set for the whole stack of words.
\end{definition}

\begin{lemma}[Agreement-to-proximity bridges]\label{lem:joint_proximity_bridges}
    \lean{ProximityGap.jointProximity_of_δ_ε_correlatedAgreementAffineLines,ProximityGap.jointProximity_of_δ_ε_correlatedAgreementCurves,ProximityGap.jointProximity_of_δ_ε_correlatedAgreementAffineSpaces}
    \uses{def:correlated_agreement}
    The correlated-agreement predicates are converted into the concrete joint-proximity witnesses
    consumed by FRI and Batched FRI soundness adapters.
\end{lemma}

\begin{theorem}[BCIKS20 Reed--Solomon proximity gap]\label{thm:bciks_rs_gap}
    \lean{ProximityGap.proximity_gap_RSCodes,ProximityGap.proximity_gap_RSCodes_johnson_of_betaRec}
    \uses{def:reed_solomon_code,def:proximity_gap,def:correlated_agreement}
    ArkLib contains the Reed--Solomon proximity-gap front doors used by STIR and Batched FRI.
    Some routes are unconditional in small-field or vacuous regimes; the full Johnson-regime curve
    correlated-agreement input is deliberately exposed as named residual hypotheses in the
    proof-system chapters.
\end{theorem}

\subsection{List-size and capacity prerequisites}\label{sec:list_size_prereqs}

\begin{theorem}[Johnson and Reed--Solomon list-size bounds]\label{thm:rs_johnson_list_size}
    \lean{ArkLib.CodingTheory.ReedSolomonJohnson.reedSolomon_johnson_list_bound,ArkLib.JohnsonBound.rs_card_ball_le,JohnsonBound.johnson_bound,JohnsonBound.card_ball_le}
    \uses{def:reed_solomon_code,def:list_decodable}
    The coding-theory tree develops Johnson-style list-size bounds and Reed--Solomon
    specializations. These results feed the proximity-gap and MCA layers rather than the
    executable protocol objects directly.
\end{theorem}

\begin{theorem}[Entropy and volume bounds]\label{thm:entropy_volume_bounds}
    \lean{CodingTheory.qEntropy,CodingTheory.choose_pow_le_qEntropy,CodingTheory.card_mul_hammingBallVolume_le_of_minDist}
    \uses{def:code_distance}
    The supporting theory includes $q$-ary entropy, Hamming-ball volume estimates, and classical
    Hamming/Plotkin/Gilbert--Varshamov style bounds. These modules support capacity and lower
    bound work around the STIR/WHIR proximity prerequisites.
\end{theorem}

\begin{remark}[Current status]\label{rem:coding_status}
    The core APIs above are active dependencies of the proof-system formalizations. The strongest
    Reed--Solomon proximity and MCA statements remain a research frontier in places, so downstream
    STIR, WHIR, and Batched FRI theorems state the missing mathematical legs as named residuals
    instead of hiding them behind untracked assumptions.  Within the identity-fold wire models the
    protocol-level soundness no longer consumes these residuals at all: the STIR checking
    verifier's round-by-round soundness is proven outright at its genuine (tight) budget and
    lifted to arbitrary security parameters by repetition
    (Theorems~\ref{thm:stir_subunit_discharge}--\ref{thm:stir_rep_wire}); the correlated-agreement
    residuals remain exactly the inputs of the degree-reducing fold regime.
\end{remark}
